\bookchapter{Глава 55. Проект 2 — приложение для погоды (сетевая версия)}

\chapterintro{Учётная книга «ставит внутреннюю силу», а приложение погоды «выводит в большой мир» — ему предстоит возиться с интернетом. В этом проекте вы научитесь получать бесплатный API, разбирать JSON, обращаться с тремя состояниями экрана «загрузка / успех / ошибка» и прятать API-ключ так, чтобы его никто не нашёл. Сделаете его — и сможете собрать любое «сетевое приложение»: новости, курсы валют, переводчик — схема везде одна.}

\booksection{Получаем бесплатный API погоды}

Для личной тренировки рекомендую OpenWeatherMap или китайский QWeather — у обоих есть бесплатные квоты: зарегистрировался и пользуешься. Пример на OpenWeatherMap:

\begin{enumerate}
  \item зарегистрируйте аккаунт на странице API keys сгенерируйте ключ (apikey);
  \item в бесплатной версии доступен метод Current Weather (текущая погода): api.openweathermap.org/data/2.5/weather? q=beijing\&appid=ВАШ\_КЛЮЧ\&units=metric;
  \item сначала попробуйте запрос в браузере и посмотрите, как выглядит ответный JSON, — обязательная «подготовка» перед подключением любого API.
\end{enumerate}

\textbf{Осторожно: ключ нельзя зашивать в код}\par

API-ключ — это «пароль от кошелька»: закоммитите его на GitHub, и за пару часов боты его вытащат и растратят. Правильно так: положите ключ в local.properties (этот файл уже исключён через gitignore) и читайте в код через buildConfig в Gradle. Найдите по запросу «лучшие практики сокрытия API-ключа в Android» и действуйте по инструкции.

\booksection{Сетевой слой: Retrofit + разбор JSON}

Сначала опишите «модель данных» по JSON, который отдаёт API, — поля, которым нет соответствия, просто игнорируйте:

\begin{lstlisting}[language=Kotlin]
// JSON: {"weather":[{"description":"Ясно"}],"main":{"temp":26.5,"humidity":40},"name":"Пекин"}
data class WeatherResponse(
    val weather: List<Weather>,
    val main: Main,
    val name: String
)
data class Weather(val description: String)
data class Main(val temp: Double, val humidity: Int)

interface WeatherApi {
    @GET("data/2.5/weather")
    suspend fun getWeather(
        @Query("q") city: String,
        @Query("appid") key: String,
        @Query("units") metric: String = "metric",
        @Query("lang") lang: String = "ru"
    ): WeatherResponse
}
\end{lstlisting}

Gson сам разложит JSON по этим data class — имена полей совпали, и магия происходит без вашего участия.

\booksection{Три состояния экрана: загрузка / успех / ошибка}

Душа сетевого приложения — «управление состояниями». Медленная сеть, нет интернета, опечатка в названии города… В каждом случае интерфейс обязан объясниться с пользователем. Изящно опишите это через sealed class (запечатанный класс):

\begin{lstlisting}[language=Kotlin]
sealed class UiState {
    object Loading : UiState()                    // крутится спиннер
    data class Success(val data: WeatherResponse) : UiState()  // успех
    data class Error(val message: String) : UiState()          // ошибка
}

// Внутри ViewModel:
fun search(city: String) {
    viewModelScope.launch {
        _state.value = UiState.Loading
        _state.value = try {
            UiState.Success(api.getWeather(city))
        } catch (e: Exception) {
            UiState.Error("Не удалось найти этот город — попробуйте другое название?")
        }
    }
}

// В интерфейсе:
when (val s = state) {
    is UiState.Loading -> CircularProgressIndicator()   // спиннер
    is UiState.Success -> WeatherContent(s.data)        // показываем
    is UiState.Error -> Text(s.message)                 // ошибка
}
\end{lstlisting}

Трёхчастная схема Loading/Success/Error — «универсальная мазь» сетевых экранов: вызубрите её, а дальше во всех сетевых приложениях делайте по этому лекалу.

\booksection{Боевой код интерфейса: рисуем погоду}

Состояния готовы; интерфейс выглядит так (ключевые фрагменты):

\begin{lstlisting}[language=Kotlin]
@Composable
fun WeatherContent(data: WeatherResponse) {
    Column(
        Modifier.fillMaxSize().padding(24.dp),
        horizontalAlignment = Alignment.CenterHorizontally
    ) {
        Text(data.name, fontSize = 28.sp, fontWeight = FontWeight.Bold)
        Text(
            "${data.main.temp}°C",
            fontSize = 72.sp,
            fontWeight = FontWeight.Light
        )
        Text(
            data.weather.firstOrNull()?.description ?: "Неизвестно",
            fontSize = 20.sp,
            color = Color.Gray
        )
        Spacer(Modifier.height(16.dp))
        Row {
            Text(" Влажность ${data.main.humidity}%")
            Spacer(Modifier.width(24.dp))
            Text(" Ощущается как ${data.main.temp}°C")
        }
    }
}
\end{lstlisting}

Несколько мелочей, у которых стоит поучиться: firstOrNull() безопасно берёт первый элемент (список может быть пуст), ?: “неизвестно” подставляет запасной текст — принцип null-безопасности действует и в интерфейсе. Температура набрана крупно, 72.sp, а «лёгкое» начертание (Light) делает цифры изящными — классический рецепт погодных приложений.

\booksection{Детали интерфейса и взаимодействия}

\booksection{Детали интерфейса и взаимодействия}

\begin{itemize}
  \item Поле поиска: TextField сверху; вводите город и жмёте «Найти» (задайте на клавиатуре действие «Поиск» — Enter сразу запускает);
  \item Карточка погоды: крупная температура (26.5°), описание погоды (ясно), влажность, название города;
  \item Избранные города: часто просматриваемые города храните в DataStore — при следующем открытии запрос уходит сам; заодно повторите главу 6;
  \item Обновление потягиванием: в Compose есть PullToRefreshBox (из Material 3) — и ощущение сразу «как в профессиональном приложении».
\end{itemize}

\begin{tcbinfo}{Совет: картинки оживляют интерфейс}
Подберите каждому типу погоды emoji или значок () и свяжите через when(description) — три строки кода, а «тепло» интерфейса удваивается. Можно подключить и стороннюю библиотеку значков погоды, но для тренировки emoji хватает.

\end{tcbinfo}

\begin{tcbpractice}{Практика: чек-лист приёмки приложения погоды}
(1) Погода нормально ищется; (2) при выключенной сети приложение изящно сообщает об ошибке и не падает (проверьте в режиме полёта!); (3) несуществующий город — дружелюбная подсказка; (4) два города в избранном остаются после перезапуска; (5) поворот экрана не теряет состояние. Пять из пяти — проект считается оконченным.

\end{tcbpractice}

\begin{tcbsummary}{Коротко о главном}
\begin{itemize}
  \item Сначала опробуйте API в браузере и разберите JSON, потом пишите код — «подготовка» ещё никого не подводила
  \item API-ключ прячьте в local.properties и ни за что не кладите в репозиторий
  \item Интерфейс Retrofit + data class: имена полей совпали — JSON разобрался сам
  \item Трёхчастное состояние Loading/Success/Error — универсальная мазь сетевых экранов
  \item Нет сети, пустой ввод, город не существует — обработка ошибок это проходной балл сетевого приложения
\end{itemize}

\end{tcbsummary}
